% Appendix skeleton. Sections marked \tbd are awaiting runs; everything else is
% final. Include from the main file after \appendix.
%
% Requires: booktabs, xcolor, algorithm/algpseudocode (shared with the main text).
\providecommand{\tbd}{\textcolor{red}{\textit{[tbd]}}}
% Citations are deliberately NOT \citep keys: no bib entries have been added.
% Each \needcite marks a reference to be looked up and entered by hand.
\providecommand{\needcite}[1]{\textcolor{red}{[cite: #1]}}

\section{Sampler details}
\label{app:smc}

\paragraph{Construction.} The triple $(q_0,\mathsf{K}_k,G_k)$ of
\S\ref{sec:latent_smc_planning} is an SMC sampler in the sense of
\needcite{Del Moral, Doucet \& Jasra, SMC samplers, JRSS-B 2006}, with $\mathsf{L}_{k-1}$ taken as the $\pi_k$-reversal of
the forward kernel. This is the standard choice: it is the tractable surrogate
for the variance-minimising backward kernel, and it is what makes the potential
computable without ever evaluating the density of $\mathsf{K}_k$
\needcite{Kim et al., arXiv:2503.15704}.

\paragraph{Invariance in practice.} The derivation assumes $\mathsf{K}_k$ leaves
$\pi_k$ invariant. Our kernel descends the planning objective and never reads
$C$, so invariance holds only approximately: the mutation moves particles
towards low objective, which is not the same as leaving $\pi_k$ fixed. We make
the assumption explicit rather than claim it, and note that the consequence ---
a mismatch between the particle distribution and $\pi_k$ --- is precisely what
the weights and the ESS-triggered resampling correct for. \tbd{} quantify the
mismatch empirically from the dumped weights.

\paragraph{Freezing the tilt.} Once $\beta_k$ reaches its cap the target sequence
has ended and further reweighting would fight the descent the kernel is
performing, so the tilt is frozen and the remaining iterations are pure
mutation. \tbd{} ablation with the freeze disabled.

\paragraph{Selection versus optimisation.} Algorithm~\ref{alg:smc} returns
$\arg\min_i \mathcal{L}(\mathbf{a}_K^{(i)})$, which is not the quantity the
sampler targets. Proposition~\ref{prop:selection} says why this is consistent:
the tilt changes which candidates are visited, never the criterion they are
judged by.

\section{Implementation details}
\label{app:impl}

The kernel block is held fixed across environments; only the particle count, the
temperature and the budget vary. Values below are the ones used for every number
in the paper.

\begin{table}[h]
  \caption{Planner settings. The kernel block is identical across environments
  by construction, so a per-environment row differs only in the roadmap graph,
  the particle count and $\alpha$.}
  \label{tab:hyper}
  \centering
  \setlength{\tabcolsep}{7pt}
  \begin{tabular}{llll}
    \toprule
    Group & Parameter & Value & Note \\
    \midrule
    \multirow{5}{*}{Kernel}
      & step size $\eta$            & $0.1$      & Adam, shared \\
      & temperature $T$             & $10^{-3}$  & Langevin noise $\sqrt{2\eta T}$ \\
      & action regulariser          & $0.1$      & \\
      & smoothing samples           & $8$ / $1$  & LeWM / DINO-WM environments \\
      & smoothing $\sigma$          & $1.0$      & annealed to $0$ \\
    \midrule
    \multirow{3}{*}{SMC}
      & ESS fraction $\tau$         & $0.5$      & resample threshold $\tau N$ \\
      & particles $N$               & $32$ / $8$ / $3$ (2 at $H{=}15$)
                                                 & Cube / Reacher / DINO-WM \\
      & temperature $\alpha$        & $0.1$ / $1$ & Cube / all others \\
    \midrule
    \multirow{4}{*}{Roadmap}
      & solver                      & heat       & \\
      & $k$-NN                      & $64$       & \\
      & graph nodes                 & $3.8$k / $12.6$k / $55$k
                                                 & Point maze / Wall / others \\
      & full-dimensional latents    & yes        & \\
    \midrule
    Budget & iterations             & $150$      & matched across all planners \\
    \bottomrule
  \end{tabular}
\end{table}

The particle count is not uniform and the non-uniformity favours the baselines:
DINO-WM retains 196 ViT patch tokens per frame through the rollout, so four
particles exceed 16\,GB of VRAM on Point maze, Wall and Push-T, while the LeWM
environments encode a single token per frame and admit 8 and 32.

\section{Baseline configurations}
\label{app:baselines}

All planners run at $150$ iterations per environment step with evaluation after
every iteration, and every arm sees the same world model, the same start and
goal pairs and the same seeds.

\paragraph{Gradient descent (GD).} Adam on the action sequence at step size
$0.1$, initialised from $\mathcal{N}(0,1)$, backpropagating through the full
$T$-step rollout. Roadmap guidance is disabled (\texttt{hjb\_tau}$=0$), so this
is the pure objective-descent baseline and differs from our method in the
sampler, not in the objective. On Cube the generic configuration diverges, so we
report a separately tuned one; the comparison there is therefore against a
strengthened baseline rather than a default.

\paragraph{GRASP.} The published configuration: joint optimisation over actions
and intermediate latent states with a Lagrangian dynamics penalty, state step
size $0.1$, action step size $10^{-3}$, and the periodic full-gradient
synchronisation its authors describe (\texttt{gd\_interval}$=100$,
\texttt{gd\_opt\_steps}$=25$).

\paragraph{CEM.} Population $30$ with $3$ elites re-fitted for $5$ outer
iterations, which is $150$ world-model rollouts per environment step and
therefore matches the budget the gradient-based planners receive. Note that the
per-iteration cost is not the same object across families: CEM spends its budget
on parallel samples and the gradient methods on sequential steps, and we match
the rollout count rather than wall-clock.

\paragraph{Budget matching.} \tbd{} whether each baseline was swept over its own
hyperparameters and the best setting reported, as is standard practice, or run
at published defaults. This must be stated explicitly: it determines whether the
comparison is against tuned or default baselines.

\section{Sampler diagnostics}
\label{app:diag}

Figure~\ref{fig:diag} reports what the sampler does per iteration on eight seeds
per environment, stratified so that the subset's success rate matches that of
the full 100-seed cell. Two arms are shown: the roadmap cost-to-go as the tilt
(ours) and the planning objective as the tilt, rescaled to the roadmap cost's
per-step mean and standard deviation so that only the \emph{ordering} of the
particles differs between them. Without that rescaling the comparison would
confound information with scale.

\begin{figure}[h]
  \centering
  \tbd
  \caption{Realised inverse temperature $\beta_k$, effective sample size against
  the $\tau N$ threshold with resample events marked, and
  $E_w[C]-E[C]$ --- the weighted-minus-uniform population mean of the roadmap
  cost --- per planner iteration. Both arms are scored on the same quantity
  $C$.}
  \label{fig:diag}
\end{figure}

\section{Ablations}
\label{app:ablations}

Unless stated otherwise each ablation runs on a 20-seed subset chosen so that
its success rate reproduces the corresponding full-table cell to within two
points (\S\ref{app:impl}); the subsets are fixed in advance and reused across
ablations.

\subsection{Roadmap versus objective}
\tbd{} Success, steps, path length and path length over straight-line distance
for the roadmap tilt, the objective tilt and no tilt at all ($\beta\equiv 0$, which
by Proposition~\ref{prop:beta0} is $N$ independent kernel chains).

\subsection{Particle count}
\tbd{} $N\in\{1,2,4,8,16,32\}$. $N{=}1$ is again the degenerate case of
Proposition~\ref{prop:beta0}: no resampling is possible, so any difference from
$N{>}1$ is attributable to the population.

\subsection{Resampling threshold, tilt strength and budget}
\label{app:knobs}

Figure~\ref{fig:abl_knobs} sweeps the three SMC knobs on Point maze $H{=}10$,
over the $20$-seed stratified subset whose success rate reproduces the full
$100$-seed cell to within two points. Each sweep moves one knob and holds the
rest at their defaults.

\paragraph{Resampling threshold.} $\tau{=}0$ never resamples, so it isolates
tempering from selection: it reaches $65\%$ against the default $\tau{=}0.5$'s
$80\%$, and $\tau{=}0.25$ is no better at $65\%$. The selection step is
therefore doing work of its own rather than merely reweighting what the kernel
already found.

\paragraph{Tilt strength.} The response is not monotone. Removing the tilt
($\beta{=}0$) costs $10$ points, confirming it matters, but the best setting is
$\beta{=}0.25$ at $95\%$, while $\beta{=}0.5$ and the default $\beta{=}1.0$ both
sit at $80\%$. The default therefore over-tilts: past a point, sharper weights
concentrate the population faster than the kernel can explore, and the extra
selection pressure buys nothing. We report the default rather than the tuned
value everywhere in this paper, so the main results understate what the method
achieves at its best setting on this cell.

\paragraph{Budget.} Success rises steeply out of a starved regime and then
flattens: $45\%$ at $25$ iterations, $80\%$ at $50$, and $85\%$ at $100$ against
the $80\%$ of the $150$ used throughout. The budget the paper runs at is thus on
the flat part of the curve rather than at a peak. We do not sweep past $150$
here, so this shows the comparison is not run at a starved budget, not that it
is insensitive to a larger one.

\begin{figure}[h]
  \centering
  \includegraphics[width=\linewidth]{figs/abl_knobs}
  \caption{SMC knobs on Point maze $H{=}10$, $20$ stratified seeds, one knob
  varied at a time. The shipped default is pink. Two of the three sweeps did not
  run their own default as an arm, so that bar comes from the main-table sweep,
  which is identical in configuration and covers these seeds; the tilt sweep
  did run its default and both sources agree at $16/20$.}
  \label{fig:abl_knobs}
\end{figure}

\section{Mutation kernel}
\label{app:kernel}

The derivation of \S\ref{sec:latent_smc_planning} holds for any $\pi_k$-invariant
$\mathsf{K}_k$, so the question is whether the result depends on which one we
use. We substitute MALA, preconditioned HMC and a learned flow transport for the
default move, holding targets, potential, schedule and resampling rule fixed.

The ablation ran on two seed strata, neither a random sample: seeds the default
kernel \emph{solves}, and seeds it \emph{fails}. We report them separately rather
than combining them. A combined estimate would credit the alternatives in a
stratum where our kernel is zero by construction --- the stratum is defined by
its failures --- so a single rescued seed out of three can read as a
double-digit deficit. Drawn apart, nothing is inferred and the asymmetry is
visible: our kernel is at $100\%$ on the left by definition and $0\%$ on the
right by definition, and only the baselines' bars carry information.

\begin{figure}[h]
  \centering
  \includegraphics[width=0.85\textwidth]{figs/fig_kernels_pointmaze} \\[3pt]
  \includegraphics[width=0.85\textwidth]{figs/fig_kernels_pusht} \\[3pt]
  \includegraphics[width=0.85\textwidth]{figs/fig_kernels_cube} \\[3pt]
  \includegraphics[width=0.55\textwidth]{figs/fig_kernels_legend}
  \caption{Success by mutation kernel, on the seeds our kernel solves (left) and
  the seeds it fails (right). Only the move differs between bars. Our kernel is a
  preconditioned, momentum-carrying discretisation of the unadjusted Langevin
  algorithm; the baselines are its Metropolis-adjusted counterpart (MALA),
  preconditioned HMC, and a learned flow transport (AFT). Strata hold three to
  ten seeds, so a single run moves a bar by $10$--$33$ points; differences of
  that size are not meaningful.}
  \label{fig:kernels}
\end{figure}

Two things survive the small samples. The Metropolis-adjusted kernels fail
\emph{completely} on Push-T --- zero of twenty runs across both strata and all
three horizons --- while reaching $80$--$90\%$ on Cube. Push-T runs two or three
particles, where a rejected proposal wastes an entire particle-iteration, and
Cube runs thirty-two, where the population absorbs rejections; the pattern
follows the particle count rather than the kernel. And on the seeds our kernel
fails, the alternatives between them rescue a substantial fraction, with AFT the
only kernel to solve anything in Push-T's failure stratum. The supportable claim
is that the construction is robust to the choice of $\pi_k$-invariant kernel and
that the kernels are \emph{complementary}, not that ours dominates.

Reacher is omitted. Its failure list predates the goal-seeding fix
(\texttt{env.action\_space.seed}), and those seeds now succeed for the default
kernel, so they are not a failure stratum.

\section{Propositions}
\label{app:proofs}

\begin{proposition}[Degenerate schedule]
\label{prop:beta0}
If $\beta_k=0$ for all $k$, Algorithm~\ref{alg:smc} reduces to $N$ independent
Markov chains driven by $\mathsf{K}_k$, with no resampling.
\end{proposition}

\begin{proof}
With $\beta_k=0$ for all $k$ we have $\Delta\beta_k=0$, so by \eqref{eq:logG}
$\log G_k=0$ and hence $G_k\equiv 1$. The weight update
\eqref{eq:smc_weight_update} then gives $W_k^{(i)}=W_{k-1}^{(i)}$, and since
$W_0^{(i)}=1$, induction yields $W_k^{(i)}=1$ for every $i$ and $k$. Normalised,
$w_k^{(i)}=1/N$, so by the definition of effective sample size
\eqref{eq:ess}, $\mathrm{ESS}_k=1/\sum_i N^{-2}=N$
\needcite{Kong, Liu \& Wong, JASA 1994 -- ESS definition}. For any threshold $\tau\in(0,1)$ we have
$\mathrm{ESS}_k=N>\tau N$, so the resampling branch never executes. Resampling
is the only step of Algorithm~\ref{alg:smc} that couples particles: every other
line acts on a single index $i$. The $N$ particles therefore evolve as
independent chains under $\mathsf{K}_k$.
\end{proof}

\begin{proposition}[The tilt selects, it does not optimise]
\label{prop:selection}
For any schedule $\{\beta_k\}$, every action sequence returned by
Algorithm~\ref{alg:smc} lies in the set reachable from the initial population by
composing $\mathsf{K}_1,\ldots,\mathsf{K}_K$, and the returned element is a
minimiser of $\mathcal{L}$ over the final population. The cost-to-go $C$ enters
neither condition.
\end{proposition}

\begin{proof}
$C$ appears in Algorithm~\ref{alg:smc} only through $C^{(i)}$, which is used to
form $\beta_k$ and $\log W_k^{(i)}$; the mutation step draws
$\bar{\mathbf{a}}_k^{(i)}\sim\mathsf{K}_k(\mathbf{a}_{k-1}^{(i)},\cdot)$, and
$\mathsf{K}_k$ is a function of the planning objective alone. Resampling
replaces the population by $\{\bar{\mathbf{a}}_k^{(A^{(i)})}\}$, a multiset of
particles already present \needcite{Douc \& Capp\'e 2005 -- resampling schemes}, and so introduces no new
points. By induction on $k$, every particle is an image of an initial draw under
a composition of the $\mathsf{K}$'s. The return step is
$\arg\min_i\mathcal{L}(\mathbf{a}_K^{(i)})$ by definition. Hence $C$ determines
which elements of the reachable set are retained, and nothing else.
\end{proof}

\section{Per-seed comparisons}
\label{app:matched}

A success rate says how often a planner solves, not \emph{which} instances it
solves. Two planners at $50\%$ may agree on every seed or disagree on all of
them, and only the second case would make the comparison a trade. Pairing by
seed separates these: Figure~\ref{fig:paired} counts, for each baseline, the
seeds we solve and it does not (W), the reverse (L), and the seeds on which the
two agree (T).

Aggregated over the grid we win $350$ and lose $70$ against gradient descent,
$444$ and $47$ against GRASP, and $344$ and $112$ against CEM. Losses are
concentrated at the longest horizon: three cells are net negative, all at
$H{=}15$ --- Point maze against CEM ($17$ vs $25$) and GRASP ($19$ vs $20$), and
Wall against CEM ($16$ vs $24$). These are the same cells in which CEM takes the
lead in Table~\ref{tab:planners}, so the paired view does not change the
picture there, it quantifies it.

\begin{figure}[h]
  \centering
  \includegraphics[width=0.32\textwidth]{figs/fig_paired_reacher_gd} \hfill
  \includegraphics[width=0.32\textwidth]{figs/fig_paired_reacher_grasp} \hfill
  \includegraphics[width=0.32\textwidth]{figs/fig_paired_reacher_cem} \\[4pt]
  \includegraphics[width=0.32\textwidth]{figs/fig_paired_pointmaze_gd} \hfill
  \includegraphics[width=0.32\textwidth]{figs/fig_paired_pointmaze_grasp} \hfill
  \includegraphics[width=0.32\textwidth]{figs/fig_paired_pointmaze_cem} \\[4pt]
  \includegraphics[width=0.32\textwidth]{figs/fig_paired_wall_gd} \hfill
  \includegraphics[width=0.32\textwidth]{figs/fig_paired_wall_grasp} \hfill
  \includegraphics[width=0.32\textwidth]{figs/fig_paired_wall_cem} \\[4pt]
  \includegraphics[width=0.32\textwidth]{figs/fig_paired_pusht_gd} \hfill
  \includegraphics[width=0.32\textwidth]{figs/fig_paired_pusht_grasp} \hfill
  \includegraphics[width=0.32\textwidth]{figs/fig_paired_pusht_cem}
  \caption{Per-seed paired outcomes, 100 seeds per cell, one panel per
  environment and baseline. Bars to the right count seeds we solve and the
  baseline does not; bars to the left count the reverse. Mirrored bars would
  mean the two methods trade instances; one-sided bars mean one solves a
  superset of the other.}
  \label{fig:paired}
\end{figure}
\section{Success against iteration, remaining panels}
\label{app:curves}

Figure~\ref{fig:success} in the main text shows one panel per environment plus
the largest remaining margin. The panels omitted there are collected here, so
the full grid is available without selecting on outcome.

\begin{figure}[h]
  \centering
  \includegraphics[width=0.315\textwidth]{figs/fig_success_reacher_h15} \hfill
  \includegraphics[width=0.315\textwidth]{figs/fig_success_pointmaze_h10} \hfill
  \includegraphics[width=0.315\textwidth]{figs/fig_success_pointmaze_h15} \\[4pt]
  \includegraphics[width=0.315\textwidth]{figs/fig_success_wall_h5} \hfill
  \includegraphics[width=0.315\textwidth]{figs/fig_success_wall_h15} \hfill
  \includegraphics[width=0.315\textwidth]{figs/fig_success_pusht_h5} \\[4pt]
  \includegraphics[width=0.315\textwidth]{figs/fig_success_pusht_h15} \\[2pt]
  \includegraphics[width=0.62\textwidth]{figs/fig_success_legend}
  \caption{Fraction of seeds solved at or before each planner iteration, for the
  environment and horizon combinations not shown in Figure~\ref{fig:success}.
  Point maze and Wall at $H{=}15$ are the two cells in which a baseline finishes
  ahead of us; in both the baseline is flat from a few iterations onward while
  our curve is still rising at the budget limit.}
  \label{fig:curves_rest}
\end{figure}
